%======================================================================
%  4. Visualization / Dashboard: design + must-have features
%======================================================================
\section{Visualization and Dashboard}
\label{sec:visualization}

\noindent The dashboard is available at \url{\dashboardurl}. The report can be
downloaded as a manual directly from the dashboard (see
Section~\ref{sec:musthaves}, feature~6).

\subsection{Design}
\label{sec:design}

TODO --- Describe how the data is displayed: the overall layout of the
dashboard, the navigation between views, the filters, and the general
ideas behind the design. For each view, justify the choice of visual
encodings (position, length, colour, size, shape) and explain why the
chosen encoding fits the question it answers.

\subsection{Must-Have Features}
\label{sec:musthaves}

\begin{table}[H]
\centering
\caption{Required features and where they are implemented.}
\label{tab:musthaves}
\begin{tabularx}{\textwidth}{@{}c X l@{}}
\toprule
\textbf{\#} & \textbf{Required feature} & \textbf{Where} \\
\midrule
1 & At least three types of graphs (e.g.\ bar chart, time series, box plot) & Figs.~\ref{fig:g1}--\ref{fig:g9} \\
2 & At least one animated graph & Fig.~\ref{fig:g7} \\
3 & An AI-generated graph, documented (platform + prompt + question answered) & Sec.~\ref{sec:aigraph} \\
4 & At least nine graphs in total & Figs.~\ref{fig:g1}--\ref{fig:g9} \\
5 & Link to the dashboard in the report & Sec.~\ref{sec:visualization} \\
6 & Option to download the report as a manual from the dashboard & TODO \\
7 & A unique, custom visualization not available as a standard R/Python function & Sec.~\ref{sec:customviz} \\
\bottomrule
\end{tabularx}
\end{table}

\subsection{The Graphs}
\label{sec:graphs}

% ---------------------------------------------------------------------
% One figure per graph. Replace the placeholder path with your image and
% keep the \label so that Table~\ref{tab:musthaves} stays correct.
% ---------------------------------------------------------------------

\begin{figure}[H]
\centering
% \includegraphics[width=\textwidth]{g1-price-timeseries.png}
\fbox{\parbox[c][4cm][c]{0.9\textwidth}{\centering TODO: insert image}}
\caption{\textbf{Graph 1 --- TODO title.} Chart type, what it encodes, and which
research question it answers.}
\label{fig:g1}
\end{figure}

\begin{figure}[H]
\centering
\fbox{\parbox[c][4cm][c]{0.9\textwidth}{\centering TODO: insert image}}
\caption{\textbf{Graph 2 --- TODO title.} TODO description.}
\label{fig:g2}
\end{figure}

\begin{figure}[H]
\centering
\fbox{\parbox[c][4cm][c]{0.9\textwidth}{\centering TODO: insert image}}
\caption{\textbf{Graph 3 --- TODO title.} TODO description.}
\label{fig:g3}
\end{figure}

\begin{figure}[H]
\centering
\fbox{\parbox[c][4cm][c]{0.9\textwidth}{\centering TODO: insert image}}
\caption{\textbf{Graph 4 --- TODO title.} TODO description.}
\label{fig:g4}
\end{figure}

\begin{figure}[H]
\centering
\fbox{\parbox[c][4cm][c]{0.9\textwidth}{\centering TODO: insert image}}
\caption{\textbf{Graph 5 --- TODO title.} TODO description.}
\label{fig:g5}
\end{figure}

\begin{figure}[H]
\centering
\fbox{\parbox[c][4cm][c]{0.9\textwidth}{\centering TODO: insert image}}
\caption{\textbf{Graph 6 --- TODO title.} TODO description.}
\label{fig:g6}
\end{figure}

% ---- Animated graph (must-have 2) -----------------------------------
\begin{figure}[H]
\centering
\fbox{\parbox[c][4cm][c]{0.9\textwidth}{\centering TODO: insert key frames}}
\caption{\textbf{Graph 7 --- animated graph.} TODO: show a few representative
frames and describe what the animation adds. State the tool used
(e.g.\ \texttt{gganimate}) and the variable mapped to time.}
\label{fig:g7}
\end{figure}

\begin{figure}[H]
\centering
\fbox{\parbox[c][4cm][c]{0.9\textwidth}{\centering TODO: insert image}}
\caption{\textbf{Graph 8 --- TODO title.} TODO description.}
\label{fig:g8}
\end{figure}

\begin{figure}[H]
\centering
\fbox{\parbox[c][4cm][c]{0.9\textwidth}{\centering TODO: insert image}}
\caption{\textbf{Graph 9 --- TODO title.} TODO description.}
\label{fig:g9}
\end{figure}

\subsection{AI-Generated Graph}
\label{sec:aigraph}

\begin{figure}[H]
\centering
\fbox{\parbox[c][4cm][c]{0.9\textwidth}{\centering TODO: insert the AI-generated graph}}
\caption{\textbf{AI-generated graph.} TODO description.}
\label{fig:aigraph}
\end{figure}

\begin{description}
  \item[Platform.] TODO --- which AI platform/tool produced the graph.
  \item[Prompt(s).] TODO --- the exact prompt(s) used, quoted verbatim.
  \item[Question answered.] TODO --- which of the questions in
        Section~\ref{sec:objectives} this graph answers.
  \item[Post-processing.] TODO --- any manual corrections applied to the output.
\end{description}

\subsection{Custom Visualization}
\label{sec:customviz}

\begin{figure}[H]
\centering
\fbox{\parbox[c][4cm][c]{0.9\textwidth}{\centering TODO: insert the custom visualization}}
\caption{\textbf{Custom visualization.} TODO description.}
\label{fig:customviz}
\end{figure}

TODO --- Explain why this visualization cannot be produced by an existing
R or Python function, how it was built (e.g.\ D3.js, hand-drawn), and what
it shows that the standard charts do not.

\subsection{Interaction}
\label{sec:interaction}

TODO --- Describe the interaction elements (filters, tooltips, brushing and
linking, zoom, year slider, region selection) with referenced screenshots.
